\chapter*{Е}
\addcontentsline{toc}{chapter}{Е}

\textbf{Еловатый}, остров. 

Еловатый остров это теперь часть Муромца около моста и въезд, начало Парка Дружбы Народов. На конец 19 века южной границей острова было современное низовье Муромца, примерно где верх улицы Трухановской. А на севером Еловатый ограничивался руслом примерно тут:

50.50252960737892, 30.53656191881135

И через русло был небольшой мост.

Западнее же остров отделялся узким руслом от широкой песчаной косы, теперь остров сросся с косой, и на месте косы ныне пляж Черторой.\\

\medskip

\textbf{Еремина, Ярёмное}, озеро-ст\'арица. 

50°23'06.1"N 30°36'26.0"E

Расположено в Осокорках, между улицей Коллекторной и 139-й Садовой. Длина озера около 480 метров. После пересечения с Коллекторной и далее на юг, озеро продолжается другой старицей, озером Пискун, а оно впадает в озеро Мартышев.